% -----------------------------------------------------------
\section{Translation, Translate}
% -----------------------------------------------------------
	\label{TRANSLATION}
	
% % ----------------------
% \subsection{See also}
% % ----------------------

% % ----------------------
% \subsection{1828 American Dictionary of the English Language} % *1828
% % ----------------------
	% \label{TRANSLATION-1828}

	% \begin{compactitem}
		% \item
	% \end{compactitem}

% % ----------------------
% \subsection{Oxford English Dictionary} % *OED
% % ----------------------
	% \label{TRANSLATION-OED}

	% \begin{compactitem}
		% \item[]
	% \end{compactitem}

% ----------------------
\subsection{Scriptural Usage} % *SCRIPTURES
% ----------------------
	\label{TRANSLATION-SCRIPTURES}

	% % -----
	% \subsubsection{Old Testament} % *OT
	% % -----
		% \label{TRANSLATION-OT}
	
		% % --
		% \paragraph{KJV}
		% % --
			% \label{TRANSLATION-OT-KJV}
	
			% \begin{tabular}{ll}
				% Total occurrences in the OT & 
			% \end{tabular}
			
			% \begin{compactdesc}
				% \item[]
			% \end{compactdesc}
			
		% % --
		% \paragraph{Hebrew Bible}
		% % --
			% \label{TRANSLATION-HEBREW}
		
			% %
			% \subparagraph{\texorpdfstring{\he{sample word}}{}}
			% %
				% \label{PAGA} % whatever is in step bible without periods
			
				% \begin{compactdesc}
					% \item[HALOT , PDF ] \textbf{}
						% \begin{compactitem}
							% \item[1]
						% \end{compactitem}
					% \item[BDB , PDF ] \textbf{}
						% \begin{compactitem}
							% \item[1]
						% \end{compactitem}
					% \item[DCH PDF ] \textbf{}
						% \begin{compactitem}
							% \item[1]
						% \end{compactitem}
				% \end{compactdesc}
				
				% \begin{compactdesc}
					% \item[Some OT uses] \textbf{}
						% \begin{compactdesc}
							% \item[]
						% \end{compactdesc}
				% \end{compactdesc}
			
		% % --
		% \paragraph{Septuagint}
		% % --
			% \label{TRANSLATION-LXX}
		
			% %
			% \subparagraph{\texorpdfstring{\gr{sample word}}{}}
			% %
		
	% % -----
	% \subsubsection{New Testament} % *NT
	% % -----
		% \label{TRANSLATION-NT}
	
		% % --
		% \paragraph{KJV}
		% % --
			% \label{TRANSLATION-NT-KJV}
			
	
			% \begin{tabular}{ll}
				% Total occurrences in the NT & 
			% \end{tabular}
			
			% \begin{compactdesc}
				% \item[]
			% \end{compactdesc}
			
		% % --
		% \paragraph{Greek New Testament} % *GREEK
		% % --
			% \label{TRANSLATION-GREEK}
		
			% %
			% \subparagraph{\texorpdfstring{\gr{sample word}}{}}
			% %
				% \label{KATALLAGE}
			
				% \begin{compactdesc}
					% \item[BDAG , PDF ] \textbf{}
						% \begin{compactitem}
							% \item 
						% \end{compactitem}
					% \item[LSJ , PDF ] \textbf{}
						% \begin{compactitem}
							% \item 
						% \end{compactitem}
				% \end{compactdesc}
				
				% \begin{compactdesc}
					% \item[Some NT uses] \textbf{}
						% \begin{compactdesc}
							% \item[]
						% \end{compactdesc}
				% \end{compactdesc}
		
	% -----
	\subsubsection{Book of Mormon} % *BOM
	% -----
		\label{TRANSLATION-BOM}
	
		% \begin{tabular}{ll}
			% Total occurrences in the BOM & 
		% \end{tabular}
		
		\begin{compactdesc}
			\item[{\hyperref[1NEPHI-13-24]{1 Nephi 13:24--28}}] \phantom{.}
				\begin{compactitem}
					\item I am wondering if this verse, or this understanding of the bible, is part of the reason that Joseph Smith undertakes the translation of the Bible later on; at the very least I would think that this verse influences that translation process. If he believes that the original Bible (whatever that means) contains the fulness of the gospel, why wouldn't he be prompted to insert it back in?
					\item It gets into verse 28 and beyond.
				\end{compactitem}
			\item[{\hyperref[MOSIAH-8-6]{Mosiah 8:6--21}}]
			\item[{\hyperref[MOSIAH-21-28]{Mosiah 21:28}}]
			\item[{\hyperref[MOSIAH-28-11]{Mosiah 28:11--20}}]
			\item[{\hyperref[ALMA-9-21]{Alma 9:21}}]
			\item[{\hyperref[MORMON-8-14]{Mormon 8:14--17}}]
			\item[{\hyperref[MORMON-9-32]{Mormon 9:32--34}}]
			\item[{\hyperref[ETHER-3-22]{Ether 3:22--4}}] \phantom{.}
				\begin{compactitem}
					\item Not everything in 4, but some stuff. 
					\item Also some stuff in 5.
				\end{compactitem}
		\end{compactdesc}
		
	% -----
	\subsubsection{Doctrine and Covenants} % *DC
	% -----
		\label{TRANSLATION-DC}
	
		% \begin{tabular}{ll}
			% Total occurrences in the DC & 
		% \end{tabular}
		
		\begin{compactdesc}
			\item[{\hyperref[DC1-29]{DC 1:29}}]
			\item[{\hyperref[DC5-4]{DC 5:4}}] 
			\item[{\hyperref[DC6-25]{DC 6:25--28}}] 
			\item[{\hyperref[DC7]{DC 7}}] 
				\begin{compactitem}
					\item This whole section is a very weird ``translation.'' See the \hyperref[DC7-HB]{historical background} and my notes. I'm guessing this ``translation'' is related to the comments in \hyperref[DC8]{DC 8} and other places, given to Cowdery, about his ``gift'' and translating ancient records. But note also that this piece doesn't get called a ``translation,'' apparently, until later. Seems to have been called a ``translation'' after the fact---it's called a ``revelation'' in the original JSP source materials.
				\end{compactitem}
			\item[{\hyperref[DC8]{DC 8}}]
				\begin{compactitem}
					\item Various verses speaking about this process.
				\end{compactitem}
			\item[{\hyperref[DC9]{DC 9}}]
				\begin{compactitem}
					\item Oliver's ``failure'' to translate; note that this whole saga of revelations involving Oliver, the gift of translation, and so on, is an important study of the translation and revelation process during this early period.
					\item Various references to ``translation'' here in this section.
				\end{compactitem}
			\item[{\hyperref[DC10]{DC 10}}]
				\begin{compactitem}
					\item More references to translation, and the ``gift.'' Note the way translation is spoken of in verses 10--11 for example.
					\item Also relevant is the ``re-translation'' and the alterations the Lord refers to.
				\end{compactitem}
			\item[{\hyperref[DC11-19]{DC 11:19, 22}}] 
			\item[{\hyperref[DC17-6]{DC 17:6}}] 
			\item[{\hyperref[DC20-8]{DC 20:8--12}}]
			\item[{\hyperref[DC21-1]{DC 21:1}}] 
			\item[{\hyperref[DC24-1]{DC 24:1}}] 
			\item[{\hyperref[DC35-20]{DC 35:20}}]
			\item[{\hyperref[DC37-1]{DC 37:1}}] 
			\item[{\hyperref[DC41-7]{DC 41:7}}] 
			\item[{\hyperref[DC42-15]{DC 42:15}}] 
			\item[{\hyperref[DC42-56]{DC 42:56--59}}]
				\begin{compactitem}
					\item The use of ``appointed'' here is interested; see \hyperref[DEATH-PHRASES]{this phrase} involving the word ``death'' and ``appoint.'' The Lord uses ``appoint'' more than once in this section and even in this passage.
				\end{compactitem}
			\item[{\hyperref[DC45-60]{DC 45:60--61}}]
			\item[{\hyperref[DC59]{DC 59}}], specifically the very end of the \hyperref[EMS-1831-JS-7-AUG-1831]{JSP version}
			\item[{\hyperref[DC73-3]{DC 73:3--4}}]
			\item[{\hyperref[DC76-15]{DC 76:15--18}}]
			\item[{\hyperref[DC77]{DC 77}}]
			\item[{\hyperref[DC84-57]{DC 84:57}}]
				\begin{compactitem}
					\item Is the Lord saying that he wrote the Book of Mormon?
				\end{compactitem} 
			\item[{\hyperref[DC90-13]{DC 90:13}}]
			\item[{\hyperref[DC91]{DC 91}}]
			\item[{\hyperref[DC94-10]{DC 94:10}}]
			\item[{\hyperref[DC107-92]{DC 107:92}}]
			\item[{\hyperref[DC124-89]{DC 124:89}}]
			\item[{\hyperref[DC128-8]{DC 128:8}}]
				\begin{compactitem}
					\item Both this use of ``translation'' and the other one in DC 128 are extremely important.
				\end{compactitem}
			\item[{\hyperref[DC128-18]{DC 128:18}}]
		\end{compactdesc}
		
	% % -----
	% \subsubsection{Pearl of Great Price} % *PGP
	% % -----
		% \label{TRANSLATION-PGP}
	
		% \begin{tabular}{ll}
			% Total occurrences in the PGP & 
		% \end{tabular}
		
		% \begin{compactdesc}
			% \item[]
		% \end{compactdesc}
		
% % ----------------------
% \subsection{Key Phrases} % *KEYPHRASES *PHRASES
% % ----------------------
	% \label{TRANSLATION-PHRASES}

	% \begin{compactdesc}
		% \item[]
	% \end{compactdesc}
	
% % ----------------------
% \subsection{Bible Dictionaries} % *DICTIONARIES
% % ----------------------
	% \label{TRANSLATION-BD}

% % ----------------------
% \subsection{Restoration Usage} % *RESTORATION
% % ----------------------
	% \label{TRANSLATION-RESTORATION}

	% % -----
	% \subsubsection{Joseph Smith} % *JS
	% % -----
		% \label{TRANSLATION-JS}
	
		% \begin{compactdesc}
			% \item[DATE/SOURCE]
		% \end{compactdesc}
	
	% % -----
	% \subsubsection{Brigham Young} % *BY
	% % -----
		% \label{TRANSLATION-BY}
	
		% \begin{compactdesc}
			% \item[DATE/SOURCE]
		% \end{compactdesc}
	
% ----------------------
\subsection{Commentary and Studies} % *COMMENTARY *STUDIES
% ----------------------
	\label{TRANSLATION-COMMENTARY}

	% -----
	\subsubsection{Key Points}
	% -----
		\label{TRANSLATION-KEYS}
	
		\begin{compactitem}
			\item
		\end{compactitem}
		
	% -----
	\subsubsection{Understandings of ``translation'' processes in the early LDS church}
	% -----
		\label{TRANSLATION-meaning}
		
		See all the DC uses above.
			
		See \hyperref[EMS-1831-JS-25-26-OCT-1831]{here}; in some minutes involving JS; also, \hyperref[EMS-1831-JS-8-NOV-1831]{here}; in a non-DC revelation \hyperref[EMS-1832-JS-CIR-20-MAR-1832]{here}; in a letter on 21 April 1833 \hyperref[EMS-1833-JS-21-APR-1833]{here}; end of section 93; JSP J1:11; see many JS journal entries in mid or late 1835---J1:107 and 109 are examples---seems pretty clear here he's using the word in the way we normally understand it. Also note that he receives some Hebrew books on 20 November, along with a Greek lexicon (this is 1835), and he's reading them at the same time as he's working with the Egyptian (see the entries for 26 and 27 November). It's pretty clear why there is Hebrew stuff in the Egyptian---he's just mixing it all together.
		
		DC 5:4, differences in older versions; there are other verses, like 30--31, in this section that also deal with this
		
		What JS says about the title page of the Book of Mormon; see \hyperref[BOMTITLE]{here}. Seems like he later in life describes the process much more as an actual translation than it really was, especially after he starts learning about other languages.
			But note places like DC 10:41, where it's clearly talking about ``engravings''; you really wish the original was extant here. But, on the other hand, the BOM itself talks a lot about engravings, and that's 1829 for sure.
			
		DC 9 is very important here too. It is specifically about the translation process, so if JS did anything along these lines, then of course the context and construction is going to be all over it from him.
		
		See JS's letter to Phelps from \hyperref[EMS-1832-JS-31-JUL-1832]{31 July 1832}.
		
		\hyperref[DC45-60]{DC 45:60--61} too. 
		
		Some verses in \hyperref[DC73]{section 73} are very explicit about the Bible work being a ``translation.''
		
		See also the statement of the translation process in \hyperref[DC76-15]{DC 76:15--18}. This is very important, as it talks about things being ``given'' through the Spirit.
		
		See \hyperref[DC24-1]{DC 24:1, 5} for comments about the Book of Mormon and about continuing to write what is given by the Comforter, which might refer back to the BOM translation process. \hyperref[DC27-5]{DC 27:5}, with ``reveal,'' is also relevant here.
		
		See \hyperref[DC90-13]{DC 90:13}, and other uses of ``translation'' listed above in the DC.
		
		See the scriptural references to ``translation'' that I have noted above.
		
		See JS's discussion of ``translation'' in his \hyperref[EMS-1840-JS-18-JUN-1840]{memo to the high council} in 1840. (See also the brief mention in the \hyperref[EMS-1840-JS-20-JUN-1840]{minutes} of a conference which occurred right after, and my notes there.)
		
		See also: \hyperref[EMS-1841-JS-16-AUG-1841-D-TS]{this mention of translation}; \hyperref[EMS-1842-JS-CIR-1-MAR-1842-DR]{this discussion}; 
		
		See the discussion (from 1838, not from 1830, as it seems) of the \hyperref[EMS-1830-JS-CIR-SPR-1830]{Book of Mormon title page}.
		
		James Burgess' account of JS, 23 July 1843, \hyperref[EMS-1843-JS-23-JUL-1843-JB]{here}, at the end (also FDR's account \hyperref[EMS-1843-JS-23-JUL-1843-FR]{here}).
		
		Ann Taves' 2016 book \textit{Revelatory Events} has a lot about translation, including lots of details and analysis on 244--247. This will be important to cover when I get to it in CtR.
		
	% -----
	\subsubsection{``Translation'' as a change of state or change of body in order to preserve bodily existence without death}
	% -----
		\label{TRANSLATION-change}
		
		Some scriptures to see:
		
			\begin{compactitem}
				\item \hyperref[ETHER-15-34]{Ether 15:34}
				\item \hyperref[DC7]{DC 7}, where John asks Christ to continue; see especially verses 2 and 6, where Christ says that he will make John ``as flaming fire.'' 
				\item \hyperref[DC107-49]{DC 107:49}
			\end{compactitem}